\documentclass[11pt]{article}
\usepackage[margin=1.1in]{geometry}
\usepackage{amsmath,amssymb,amsthm}
\usepackage[round]{natbib}
\usepackage{hyperref}
\hypersetup{hidelinks}
\usepackage{microtype}

\newtheorem{theorem}{Theorem}
\newtheorem{proposition}[theorem]{Proposition}
\newtheorem{lemma}[theorem]{Lemma}
\newtheorem{corollary}[theorem]{Corollary}
\theoremstyle{definition}
\newtheorem{definition}[theorem]{Definition}
\newtheorem{assumption}{Assumption}
\theoremstyle{remark}
\newtheorem{remark}[theorem]{Remark}

\newcommand{\E}{\mathbb{E}}
\newcommand{\Prob}{\mathbb{P}}
\newcommand{\Var}{\mathrm{Var}}
\newcommand{\Cov}{\mathrm{Cov}}
\newcommand{\one}{\mathbf{1}}
\newcommand{\given}{\,|\,}
\newcommand{\Fcal}{\mathcal{F}}

\title{Conformal Prediction as an Averaged Limit:\\
notes on the corrector under Markov dependence}
\author{Peter Cotton\\\texttt{peter.cotton@microprediction.com}}
\date{September 2026}

\begin{document}
\maketitle

\begin{abstract}
Pooling calibration residuals averages a fast variable over its invariant law, which is the
averaged limit of homogenization. Split conformal prediction is therefore correct to leading order
when the fast variable is fast relative to the forecast horizon, and the gap it leaves is a
corrector. For a Markov chain that corrector solves a Poisson equation, and the same solution
generates the calibration variance. On GARCH at the persistence fitted to daily equities the
marginal guarantee holds exactly while conditional coverage runs from 98\% to 74\% across
volatility deciles. Four attempts to convert the corrector into a calibration procedure are
recorded, with what defeated each.
\end{abstract}

\section{Relation to existing work}\label{sec:status}

Three of the quantities below are classical and we claim none of them. Effective sample size for
autocorrelated data is \citet{bayley1946effective}, in the title. The mean return time to a set is
Kac's lemma. The clustering of visits to a set is measured by the index of dispersion of counts.

The conformal application is also in print. \citet{ramos2024transported} study the
calibration-conditional coverage of a conformal threshold under stationary mixing, with the
long-run indicator variance, clustered calibration, an effective sample size and Berry--Esseen
approximations. Anyone wanting a usable dependence correction for conformal calibration should
start there.

What may be unclaimed is narrow. The exact Poisson representation of realized history-conditional
coverage in Theorem~\ref{thm:corrector} does not appear in the conformal literature, nor does the
observation that the same solution family supplies the calibration variance. The Markov chain
literature has not been searched for the identity itself, so it should be assumed to exist.

\section{The exact object}

Take a stationary time series of nonconformity scores and the split-conformal threshold, an
order statistic of a calibration block. Two things are true at once and are usually
conflated. The calibration scores are dependent, so the order statistic used as the threshold
is not the one an i.i.d.\ analysis assumes. And the test point is generated by a dynamical
process whose conditional law given the current state need not equal the invariant mixture.
The exact object of study is neither the invariant law nor the marginal coverage. It is the
realized history-conditional coverage of a random order-statistic threshold, and both effects
act on it.

It is tempting to say that ordinary conformal prediction ``calibrates the invariant law.''
Under dependence that is not an exact guarantee. Dependence alters the law of the order
statistic, and the conditional law given the current state alters the law of the future test
score. Our contribution is that one solution family of a Poisson equation governs both
perturbations: the transient gap between the current-state and invariant score laws, and the
long-run calibration noise of the dependent empirical quantile. The identity is exact and
asymptotics-free, and we verify every algebraic claim numerically to machine precision.

\paragraph{Neighbouring work.} Distributional conformal prediction for Markov processes builds a
conditional probability integral transform from estimated transition CDFs \citep{dai2024markov}, and
concentration-based penalties quantify coverage loss under non-exchangeable and Markovian data
\citep{oliveira2024split,zheng2024markovian}. The nonparametric price of a fixed-point guarantee is
the subject of a companion note \citep{cotton2026honest}, whose local modulus reappears in
Section~\ref{sec:modulus}. Section~\ref{sec:status} places these notes against
\citet{ramos2024transported}, which is the nearest work.

\section{The exact coverage decomposition}\label{sec:decomp}

Let $(W_t)$ be a Markov chain with transition kernel $P$ and invariant distribution $\pi$,
on a state space rich enough that all emission randomness is absorbed, so the score is a
deterministic readout $R_t=r(W_t)$. This convention matters below: with the emission
randomness inside $W_t$, the calibration-variance calculation needs no separate emission
martingale. For each threshold $q$ write the coverage function and the invariant score CDF
\begin{equation}\label{eq:Hdef}
  H_q(w) = \one\{r(w)\le q\},
  \qquad
  F(q) = \pi H_q .
\end{equation}

\begin{assumption}[Ergodic chain with bounded Poisson solutions]\label{as:chain}
The chain is uniformly ergodic with invariant law $\pi$, and there is a jointly measurable map
$(w,q)\mapsto\chi_q(w)$ solving the Poisson equation
\begin{equation}\label{eq:poisson}
  (I-P)\chi_q = H_q - F(q), \qquad \pi\chi_q = 0 ,
\end{equation}
with $\sup_q\|\chi_q\|_\infty<\infty$ over the quantile range used. For a finite irreducible
aperiodic chain both hold automatically.
\end{assumption}

Let $R_{(k)}$ be the $k$th order statistic of the calibration scores $R_1,\dots,R_n$, let
$\Fcal_n$ be the calibration history, and let $\mu_n(\cdot)=\Prob(W_n\in\cdot\given\Fcal_n)$
be the filtering distribution of the current state. If the test score is observed $\ell$
steps after the calibration block, its realized coverage is
\begin{equation}\label{eq:Cdef}
  C_{n,\ell,k} = \Prob\{R_{n+\ell}\le R_{(k)}\given \Fcal_n\}
    = \mu_n\big(P^{\ell}H_{R_{(k)}}\big),
\end{equation}
the $\ell$-step transported coverage function averaged against the filter. The invariant law
is the wrong reference by exactly the solved Poisson equation. Since $F(q)$ is constant and
$P^\ell$ fixes constants,
\begin{equation}\label{eq:disc}
  P^{\ell}H_q - F(q) = P^{\ell}\big(H_q-F(q)\big) = (P^{\ell}-P^{\ell+1})\chi_q .
\end{equation}
The first equality already writes the state discrepancy without solving anything. What the
Poisson solution $\chi_q$ adds, and the reason we carry it, is that the \emph{same} $\chi_q$
generates the additive-functional martingale and its variance in
Section~\ref{sec:onepoisson}.

\begin{theorem}[Exact conformal coverage corrector]\label{thm:corrector}
Under Assumption~\ref{as:chain}, for every calibration threshold $R_{(k)}$,
\begin{equation}\label{eq:corrector}
  C_{n,\ell,k} = F(R_{(k)}) + \mu_n\big[(P^{\ell}-P^{\ell+1})\chi_{R_{(k)}}\big] .
\end{equation}
The first term is the coverage of the random calibration threshold under the invariant score
law. The second is the exact current-state correction. If the state is observed at time $n$,
so $W_n$ is $\Fcal_n$-measurable, the filter is a point mass and
\begin{equation}\label{eq:corrector-obs}
  C_{n,\ell,k} = F(R_{(k)}) + (P^{\ell}-P^{\ell+1})\chi_{R_{(k)}}(W_n) .
\end{equation}
\end{theorem}

No asymptotics enter. Given the Poisson solution, \eqref{eq:corrector} is an algebraic
identity holding for every $n$, $\ell$, $k$. A direct simulation on a finite chain confirms
both sides of \eqref{eq:corrector-obs} agree to machine precision
(Section~\ref{sec:repro}).

\section{A transported order-statistic representation}\label{sec:transport}

Suppose $F$ is continuous and set $U_i=F(R_i)$. Each $U_i$ is marginally uniform under
stationarity, though the sequence need not be independent, and monotonicity of $F$ gives
$F(R_{(k)})=U_{(k)}$. Define the state-dependent coverage transport
\begin{equation}\label{eq:tau}
  \tau_{\ell,w}(u) = P^{\ell}H_{F^{-1}(u)}(w)
    = u + (P^{\ell}-P^{\ell+1})\chi_{F^{-1}(u)}(w) .
\end{equation}
Then \eqref{eq:corrector-obs} reads $C_{n,\ell,k}=\tau_{\ell,W_n}(U_{(k)})$, which separates
the problem cleanly. The law of $U_{(k)}$ carries the calibration dependence. The map
$\tau_{\ell,W_n}$ carries the test-side state dependence. Under i.i.d.\ calibration $U_{(k)}$
has the familiar beta law and $\tau$ is the identity, recovering ordinary conformal coverage.
Under dependent calibration the law of $U_{(k)}$ is generally not beta, and its leading
variance is set by the long-run variance of the coverage indicator, computed in
Section~\ref{sec:onepoisson}. This is the transported-beta separation
\citep{ramos2024transported}, to which the Poisson representation adds the explicit corrector
and a martingale representation of the calibration term.

\section{The first-order coverage expansion}\label{sec:expansion}

Fix $p=1-\alpha$ and $q_p=F^{-1}(p)$, and abbreviate the transient corrector
\begin{equation}\label{eq:bias}
  b_{\ell,w}(q) = P^{\ell}H_q(w) - F(q) = (P^{\ell}-P^{\ell+1})\chi_q(w) .
\end{equation}
The dependent empirical quantile obeys a Bahadur representation
\begin{equation}\label{eq:bahadur}
  F(R_{(k_p)}) - p = -\big[F_n(q_p) - p\big] + o_p(n^{-1/2}),
  \qquad F_n(q) = \frac1n\sum_{i=1}^n \one\{R_i\le q\},
\end{equation}
valid under geometric ergodicity with $F$ differentiable at $q_p$ and $f(q_p)>0$
\citep{glynn2002quantile}. The expansion below differentiates the corrector in the threshold, so
it needs more than \eqref{eq:bahadur} supplies. We assume in addition that, in a neighbourhood of
$q_p$, the map $q\mapsto b_{\ell,w}(q)$ is continuously differentiable uniformly in $w$, that $F$ is
continuously differentiable with $f$ bounded away from zero there, and that
$R_{(k_p)}\to q_p$ in probability, which the same geometric ergodicity gives. For a finite chain
the first condition reduces to continuous differentiability of each $F_w$ near $q_p$, since
$b_{\ell,w}$ is then a finite linear combination of the $F_{w'}$. Under these, expanding
\eqref{eq:corrector-obs} around $q_p$ and using \eqref{eq:bahadur},
\begin{equation}\label{eq:expansion}
  C_{n,\ell,k_p} - p
  = b_{\ell,W_n}(q_p)
    - \Big[1 + \tfrac{1}{f(q_p)}\partial_q b_{\ell,W_n}(q_p)\Big]\big[F_n(q_p)-p\big]
    + o_p(n^{-1/2}).
\end{equation}
When the transient correction is small in $C^1$, $\|b_{\ell,\cdot}\|_{C^1}=o(1)$, this reduces
to the central formula
\begin{equation}\label{eq:central}
  C_{n,\ell,k_p} - p = b_{\ell,W_n}(q_p) - \big[F_n(q_p)-p\big]
    + o_p\big(\|b_{\ell,\cdot}\| + n^{-1/2}\big) .
\end{equation}
Realized coverage error has two first-order sources, a current-state bias $b_{\ell,W_n}(q_p)$
and a calibration-CDF error $F_n(q_p)-p$. The first-order probability-level correction is
$p\mapsto p-b_{\ell,W_n}(q_p)$. If the current state causes undercoverage then
$b_{\ell,W_n}(q_p)<0$ and the required empirical quantile level rises. A one-sided honesty
margin is then added on top, and its size is set by the same operator.

\section{One Poisson equation}\label{sec:onepoisson}

No second Poisson equation is required. Center the coverage indicator at the invariant
quantile, $g_p(w)=H_{q_p}(w)-p$. Because $F(q_p)=p$, equation \eqref{eq:poisson} at $q=q_p$
reads $(I-P)\chi_{q_p}=H_{q_p}-p=g_p$, so by uniqueness of the centered solution the
calibration-variance potential is the corrector itself,
\begin{equation}\label{eq:psi-is-chi}
  \psi_p = \chi_{q_p} .
\end{equation}
The single solution family $q\mapsto\chi_q$ therefore supplies both objects,
\begin{equation}\label{eq:unified}
  b_{\ell,w}(q) = (P^{\ell}-P^{\ell+1})\chi_q(w),
  \qquad
  \sigma^2(q) = \pi\big[\chi_q^2 - (P\chi_q)^2\big] .
\end{equation}
The additive calibration error has the exact martingale decomposition, indexed to match the
calibration block $i=1,\dots,n$,
\begin{equation}\label{eq:martingale}
  \sum_{i=1}^{n} g_p(W_i) = P\chi_{q_p}(W_0) - P\chi_{q_p}(W_n) + \sum_{i=1}^{n} D_i,
  \qquad D_i = \chi_{q_p}(W_i) - P\chi_{q_p}(W_{i-1}),
\end{equation}
where $(D_i)$ are martingale differences for the natural filtration and the endpoint term is
bounded by $2\|\chi_{q_p}\|_\infty$. The long-run variance is
\begin{equation}\label{eq:lrv}
  \sigma_p^2 = \sigma^2(q_p)
    = p(1-p) + 2\sum_{j\ge1}\Cov_\pi\big(\one\{R_0\le q_p\},\one\{R_j\le q_p\}\big) .
\end{equation}
So one operator supplies the transient bias $b_{\ell,w}$ and the calibration variance
$\sigma_p^2$. These are not unrelated corrections. The corrector measures how much fast-state
information survives to the test horizon. The variance measures how much fast-state
persistence removes calibration information. For a reversible chain they are the same
coefficients read two ways.

\begin{proposition}[Reversible spectral reading]\label{prop:spectral}
Let $P$ be reversible with respect to $\pi$ and suppose its spectrum is pure point, with real
eigenvalues $1=\lambda_1>\lambda_2\ge\cdots$ and a $\pi$-orthonormal eigenbasis $\phi_j$ of
$L^2_0(\pi)$. A finite irreducible aperiodic chain, or any $P$ with compact resolvent, satisfies
this. For a chain with continuous spectrum, replace the sums below by integrals against the
spectral measure of $g_q$, which leaves both identities intact in that form. Expand the centered coverage function
$g_q=H_q-F(q)=\sum_{j\ge2}c_j(q)\phi_j$. Then
\begin{equation}\label{eq:spectral}
  b_{\ell,w}(q) = \sum_{j\ge2} c_j(q)\,\lambda_j^{\ell}\,\phi_j(w),
  \qquad
  \sigma^2(q) = \sum_{j\ge2} c_j(q)^2\,\frac{1+\lambda_j}{1-\lambda_j} .
\end{equation}
\end{proposition}

\begin{proof}
Solving \eqref{eq:poisson} gives $\chi_q=\sum_{j\ge2}c_j(q)(1-\lambda_j)^{-1}\phi_j$, so
$(P^\ell-P^{\ell+1})\chi_q=\sum_j c_j(q)\lambda_j^\ell(1-\lambda_j)(1-\lambda_j)^{-1}\phi_j$,
the first identity. For the second, $\pi[\chi_q^2-(P\chi_q)^2]=\sum_j
c_j(q)^2(1-\lambda_j^2)(1-\lambda_j)^{-2}=\sum_j c_j(q)^2(1+\lambda_j)/(1-\lambda_j)$.
\end{proof}

The same coefficients $c_j(q)$ and eigenvalues $\lambda_j$ control the decay of current-state
memory through $\lambda_j^\ell$ and the inflation of calibration variance through
$(1+\lambda_j)/(1-\lambda_j)$. The slogan that state bias and effective-sample-size loss are
two readings of one fast mode is exact when a single non-unit eigenvalue dominates, as in the
two-state model below, and otherwise they are two functionals of the same spectrum.

\section{A fast two-state volatility model}\label{sec:twostate}

The cleanest one-mode instance is a fast volatility regime. Let $V^\varepsilon_t\in\{L,H\}$
be a continuous-time chain with generator
\begin{equation}\label{eq:generator}
  Q_\varepsilon = \frac1\varepsilon\begin{pmatrix} -a & a\\ b & -b\end{pmatrix},
  \qquad \kappa = a+b, \quad \pi_H = \frac{a}{a+b},\quad \pi_L=\frac{b}{a+b}.
\end{equation}
Conditional on the regime the score CDFs are $F_L(q)$ and $F_H(q)$, and the invariant score
law is the mixture $F(q)=\pi_L F_L(q)+\pi_H F_H(q)$. We assume emissions are conditionally
independent across time given the regime sequence, or equivalently that any remaining
dependence has been folded into the augmented Markov state, so that \eqref{eq:cov} below
holds. If the filtered high-state probability at time zero is
$\omega_0=\Prob(V_0=H\given\Fcal_0)$, then at horizon $t$
\begin{equation}\label{eq:omega}
  \omega_t = \pi_H + e^{-\kappa t/\varepsilon}(\omega_0 - \pi_H) ,
\end{equation}
so the exact conditional score CDF and transient coverage error are
\begin{equation}\label{eq:Gt}
  G^\varepsilon_t(q) = F(q) + e^{-\kappa t/\varepsilon}(\omega_0-\pi_H)\big[F_H(q)-F_L(q)\big],
  \qquad
  b^\varepsilon_t(q) = G^\varepsilon_t(q) - F(q) .
\end{equation}
This is a literal homogenization boundary layer. After several fast mixing times it is
exponentially small. At a forecast horizon comparable to $\varepsilon$ it is not small at
all. Where $F$ is differentiable with $f>0$ near $q_p$, and treating $e^{-\kappa t/\varepsilon}$
as small (that is, after enough fast mixing, not uniformly at $t=0$), the quantile correction
is
\begin{equation}\label{eq:qcorr}
  q^\varepsilon_{t,p} = q_p - \frac{b^\varepsilon_t(q_p)}{f(q_p)} + O\big(e^{-2\kappa t/\varepsilon}\big),
  \qquad f = \pi_L f_L + \pi_H f_H .
\end{equation}
For a nonnegative absolute-error score the high-volatility state usually has $F_H(q)<F_L(q)$
at the relevant $q$. If $\omega_0>\pi_H$ then $b^\varepsilon_t(q_p)<0$ and \eqref{eq:qcorr}
raises the threshold, as it should.

\paragraph{Discrete sampling and calibration dependence.} Sample every $\Delta$ time units.
The nontrivial eigenvalue of the sampled transition matrix is $\lambda=e^{-\kappa\Delta/
\varepsilon}$. Writing $d_p=F_H(q_p)-F_L(q_p)$, the coverage indicator
$I_j=\one\{R_j\le q_p\}$ has
\begin{equation}\label{eq:cov}
  \Cov(I_0,I_j) = \pi_L\pi_H d_p^2\,\lambda^{j},\quad j\ge1,
  \qquad
  \sigma_p^2 = p(1-p) + 2\pi_L\pi_H d_p^2\,\frac{\lambda}{1-\lambda} .
\end{equation}
One fast eigenmode produces both effects, the transient bias $\propto\lambda^{m}(\omega_0-
\pi_H)d_p$ and the variance inflation $\propto\pi_L\pi_H d_p^2\lambda/(1-\lambda)$, in
agreement with Proposition~\ref{prop:spectral} at a single mode. When $\Delta\ll\varepsilon$
we have $1-\lambda\sim\kappa\Delta/\varepsilon$, so $\lambda/(1-\lambda)\sim\varepsilon/
(\kappa\Delta)$. Sampling faster piles up redundant calibration scores, and for a window of
fixed physical duration $T=n\Delta$ the effective sample size saturates at order
$T/\varepsilon$ rather than growing like $n$.

\paragraph{A numerical illustration.} Take $p=0.9$, $\pi_H=0.2$, $\lambda=0.8$, and at the
invariant $0.9$-quantile $F_L(q_p)=0.95$, $F_H(q_p)=0.70$, so the mixture is
$0.8(0.95)+0.2(0.70)=0.90$. If the current state is high volatility the one-step coverage of
the invariant threshold is $0.9+0.8(1-0.2)(0.70-0.95)=0.74$. If it is low volatility it is
$0.9+0.8(0-0.2)(0.70-0.95)=0.94$. The same pooled threshold gives either $74\%$ or $94\%$
conditional coverage while appearing perfectly calibrated at $90\%$ in stationarity. The
long-run variance is $\sigma_p^2=0.09+2(0.8)(0.2)(0.25)^2(0.8/0.2)=0.17$, an integrated
dependence factor $\sigma_p^2/[p(1-p)]\approx1.89$, so the nominal calibration sample behaves
like only $\approx53\%$ as many independent observations. A direct simulation of the sampled
chain reproduces $\sigma_p^2$ (Section~\ref{sec:repro}). The persistent regime creates severe
current-state undercoverage and nearly halves the effective calibration sample at once.

\section{Homogenization: the continuous-time limit}\label{sec:homog}

The two-state calculation is the visible corner of a general fast--slow limit. Let $Q$ be the
generator of an ergodic continuous-time chain with invariant law $\pi$, sampled every
$\Delta$ at fast scale $\varepsilon$, so the sampled kernel is
\begin{equation}\label{eq:sampled}
  P_{\Delta,\varepsilon} = e^{(\Delta/\varepsilon)Q},
  \qquad \delta = \Delta/\varepsilon .
\end{equation}
For the centered coverage function $g_q=H_q-F(q)$ define the continuous-time cell problem
\begin{equation}\label{eq:cell}
  -Q\,\varphi_q = g_q, \qquad \pi\varphi_q = 0 .
\end{equation}

\begin{proposition}[Poisson-to-cell convergence]\label{prop:homog}
On the centered subspace, as $\delta\downarrow0$,
\begin{equation}\label{eq:homlim}
  \delta\,\chi_q \;\longrightarrow\; (-Q)^{-1}g_q = \varphi_q ,
\end{equation}
where $\chi_q$ solves \eqref{eq:poisson} for $P_{\Delta,\varepsilon}$. If $Q$ is reversible,
the per-observation long-run variance obeys $\delta\,\sigma^2_\delta(q)\to
2\langle g_q,(-Q)^{-1}g_q\rangle_\pi$, so for a calibration window of physical duration
$T=n\Delta$,
\begin{equation}\label{eq:greenkubo}
  \Var\big(F_n(q)\big) \approx \frac{2\varepsilon}{T}\,\langle g_q,(-Q)^{-1}g_q\rangle_\pi .
\end{equation}
At physical test horizon $t=m\Delta$ the transient corrector is the fast boundary layer
$b^\varepsilon_t(q,w) = e^{tQ/\varepsilon}g_q(w)$.
\end{proposition}

\begin{proof}
On the centered subspace $I-e^{\delta Q}=-\delta Q(I+O(\delta))$, so
$\delta(I-e^{\delta Q})^{-1}\to(-Q)^{-1}$ and $\delta\chi_q=\delta(I-e^{\delta Q})^{-1}g_q
\to(-Q)^{-1}g_q$. For reversible $Q$, $\sigma^2_\delta(q)=\langle g_q,g_q\rangle_\pi+2\sum_{j
\ge1}\langle g_q,e^{j\delta Q}g_q\rangle_\pi$, a Riemann sum for $\delta^{-1}\int_0^\infty
2\langle g_q,e^{sQ}g_q\rangle_\pi\,ds=\delta^{-1}\cdot2\langle g_q,(-Q)^{-1}g_q\rangle_\pi$.
Dividing by $n=T/\Delta$ gives \eqref{eq:greenkubo}. The boundary layer is
$P^m_{\Delta,\varepsilon}g_q=e^{m\Delta Q/\varepsilon}g_q$.
\end{proof}

Three readings follow. The discrete Poisson solution converges to the cell solution. The
honesty variance converges to the Green--Kubo coefficient. The current-state correction is
the fast boundary layer. The two-state model of Section~\ref{sec:twostate} is the one-mode
corollary, and \eqref{eq:greenkubo} is the general form of the effective-sample-size
saturation at order $T/\varepsilon$. A slow state $X^\varepsilon_t$ can be added through a
frozen-$x$ generator $Q_x$ and its cell problem, with a remainder controlling the drift of
$x$ over the calibration and test horizons. We do not carry that step here.

\section{GARCH: the corrector made visible}\label{sec:garch}

The two-state model of Section~\ref{sec:twostate} is a cartoon. The reading it illustrates is
easiest to see on a process practitioners actually fit. Take GARCH(1,1),
\[
  r_t = \sigma_t \varepsilon_t, \qquad
  \sigma_t^2 = \omega + \alpha r_{t-1}^2 + \beta \sigma_{t-1}^2,
  \qquad \omega = 1-\alpha-\beta ,
\]
so the stationary variance is one, with score $R_t = |r_t|$ and a pooled split-conformal threshold
at level $p=0.9$ from $n=2000$ calibration points. GARCH is convenient here because $\sigma_{t+1}$
is measurable at time $t$, so the conditional coverage of a threshold $T$ is exactly
$2\Phi(T/\sigma_{t+1})-1$ with no estimation anywhere. It is also the adversarial case for the
averaged reading, because daily equity persistence $\alpha+\beta \approx 0.99$ gives a relaxation
half-life near seventy days against a one-day horizon, which is deep inside the boundary layer.

\begin{center}
\begin{tabular}{r r r r r r}
\hline
$\alpha+\beta$ & half-life & calmest decile & median decile & most volatile decile & marginal \\
\hline
0.800 & 3.1   & 0.92 & 0.91 & 0.86 & 0.900 \\
0.900 & 6.6   & 0.93 & 0.91 & 0.84 & 0.899 \\
0.960 & 17.0  & 0.96 & 0.91 & 0.80 & 0.900 \\
0.990 & 69.0  & 0.98 & 0.92 & 0.74 & 0.904 \\
0.997 & 230.7 & 0.99 & 0.93 & 0.71 & 0.901 \\
0.999 & 692.8 & 0.99 & 0.93 & 0.72 & 0.896 \\
\hline
\end{tabular}
\end{center}

Deciles are of the current volatility state, nominal coverage is $0.90$, and the last column is the
marginal rate. The marginal guarantee holds exactly everywhere, so conformal
does precisely what it promises and no row is evidence of a bug. And the spread across a row is the
corrector of Theorem~\ref{thm:corrector}: it vanishes as the volatility mixes faster, and at the
persistence actually fitted to daily equities the same interval delivers $98\%$ coverage in calm
markets and $74\%$ in volatile ones.

Pooling averages the fast variable, averaging is valid when
the fast variable is fast, and the error is the corrector. The spread grows monotonically with the
relaxation time and then saturates once the half-life approaches the calibration window, where the
sample no longer sees the stationary law. We swept persistence by moving both GARCH parameters,
which also moves the stationary volatility-of-volatility, so the apparent square-root scaling of the
spread in the half-life should not be quoted without a design that holds the latter fixed.

\section{A finite-sample honest Markov correction}\label{sec:honest}

The expansion motivates a finite-sample theorem. We ask for protection at one invariant
probability level, not uniformly over every threshold, which is both sharper and cleaner.

\begin{lemma}[One-sided Freedman corollary]\label{lem:freedman}
Let $(D_i)_{i\le n}$ be martingale differences with $|D_i|\le c$ and predictable quadratic
variation $\sum_i\E[D_i^2\given\Fcal_{i-1}]\le nV$. Then for $x=\log(1/\eta)$, with
probability at least $1-\eta$, $\sum_{i=1}^n D_i \le \sqrt{2nVx}+\tfrac{c}{3}x$.
\end{lemma}

Write $x=\log(1/\eta)$ and, with $B=\|\chi_{q}\|_\infty$ and $V$ a uniform bound on the
conditional increment variance (specified per model below),
\begin{equation}\label{eq:delta}
  \delta_{n,\eta} = \frac{2B + \sqrt{2nVx}}{n} + \frac{2Bx}{3n} .
\end{equation}
For a model class $\Theta$ of dependent score dynamics define the target-specific worst-case
level
\begin{equation}\label{eq:uwc}
  u^\star = \inf\Big\{ u\in[p,1) :\;
    \inf_{\theta\in\Theta}\ \inf_{w}\ P_\theta^{\ell}H_{\theta,F_\theta^{-1}(u)}(w) \ge p \Big\},
\end{equation}
the smallest invariant level such that, whatever the model and current state, the $\ell$-step
coverage of the corresponding invariant quantile is at least $p$.

\begin{theorem}[Finite-sample class-uniform honest correction]\label{thm:honest}
Suppose each $F_\theta$ is continuous and strictly increasing on the relevant range, that
$u^\star+\delta_{n,\eta}<1$, and that the increments $D_i=\chi_{\theta,F_\theta^{-1}(u^\star)}
(W_i)-P_\theta\chi_{\theta,F_\theta^{-1}(u^\star)}(W_{i-1})$ obey $|D_i|\le 2B$ and conditional
variance at most $V$, uniformly over $\theta\in\Theta$. Set $k=\lfloor n(u^\star+
\delta_{n,\eta})\rfloor+1$, with $R_{(k)}=+\infty$ if $k>n$. Then
\begin{equation}\label{eq:honest-guarantee}
  \inf_{\theta\in\Theta}\ \Prob_\theta\Big\{\, \Prob_\theta\big(R_{n+\ell}\le R_{(k)}\given
    \Fcal_n\big) \ge p \,\Big\} \ge 1-\eta .
\end{equation}
The procedure uses only the rank $k$, not the model.
\end{theorem}

\begin{proof}
Fix $\theta$ and write $u=u^\star$, $q_u=F_\theta^{-1}(u)$, which is finite because $u<1$ and
$F_\theta$ is continuous and strictly increasing, and $F_\theta(q_u)=u$ so $g_u=H_{q_u}-u$ is
centered. By \eqref{eq:martingale} and Lemma~\ref{lem:freedman}, with probability at least
$1-\eta$,
\[
  F_n(q_u) \le u + \delta_{n,\eta} .
\]
On this event $nF_n(q_u)\le n(u+\delta_{n,\eta})<k$, so fewer than $k$ calibration scores lie
at or below $q_u$, hence $R_{(k)}\ge q_u=F_\theta^{-1}(u^\star)$. Since $q\mapsto P_\theta^\ell
H_q(w)$ is nondecreasing,
\[
  \Prob_\theta\big(R_{n+\ell}\le R_{(k)}\given\Fcal_n\big)
    = P_\theta^{\ell}H_{R_{(k)}}(W_n)
    \ge P_\theta^{\ell}H_{F_\theta^{-1}(u^\star)}(W_n)
    \ge p ,
\]
the last step by the definition \eqref{eq:uwc} of $u^\star$. The bound is uniform in
$\theta$.
\end{proof}

\begin{corollary}[Uniform-in-threshold form]\label{cor:al}
With $a_\ell=\sup_{w,q}[F(q)-P^\ell H_q(w)]_+$ and $u=p+a_\ell<1$, the choice $u^\star\le
p+a_\ell$ recovers the simpler calibration level $p+a_\ell$, since $P^\ell H_{F^{-1}(p+a_\ell)}
(w)\ge F(F^{-1}(p+a_\ell))-a_\ell=p$.
\end{corollary}

This is a genuine calibration-sample honest and state-conditionally valid guarantee, stronger
than a statement about average marginal coverage, and uniform over the class. One caution.
The procedure is honest because $k$ depends only on the fixed class constants $u^\star,B,V$.
If instead $P$, $F$, and the Poisson solution are estimated from the same data, finite-sample
validity is not automatic. It then requires deterministic uniform class bounds, a
simultaneous confidence set for the dynamics, or a separate estimation block with a valid
envelope.

\section{Four attempts that failed}\label{sec:failed}

Theorem~\ref{thm:honest} is correct, but its levels are class constants computed from a known
kernel, so it is not a procedure anyone can run. Four attempts to turn the corrector into one are
recorded here, with what killed each. The scripts are in the repository and
\texttt{EXPERIMENTS.md} carries the full numbers.

\paragraph{Index the calibration level by the observed state.} Keep a pooled calibration block
and choose, for each state, the smallest level whose invariant quantile is $\ell$-step adequate from
that state, paying a union bound over states inside \eqref{eq:delta}. This is honest, by
Theorem~\ref{thm:honest} applied simultaneously at $M$ fixed thresholds. It is also dominated by a
simpler scheme that needs no Poisson solution at all: pair each calibration score with the state
$\ell$ steps earlier and take an ordinary quantile within each lagged-state stratum. The indexed
scheme can only slide along the quantile function of the \emph{invariant} law, so in a state that
already over-covers its level cannot fall below $p$ and the construction stops, while stratifying
estimates the $\ell$-step law directly and continues. In the two-state model at $\pi_H=0.2$ the
indexed scheme has expected threshold $1.155$ against the stratified scheme's $1.002$ and
over-covers at $0.940$ where stratifying is exact at $0.900$. No sample size we tried reverses it.

\paragraph{Claim the sample-size advantage over stratifying.} Stratifying splits the block, so a
stratum holds about $n\pi(w)$ points, while a level-indexed pooled scheme uses all $n$ and pays only
$\sqrt{\log M}$ in the deviation. This is not a statistical gain. The pooled scheme is handed the
per-state levels as population objects computed from a known kernel, so it never pays to learn $M$
quantiles. Minimax lower bounds for conditional conformal coverage over a class of dimension $d$
scale as $\sqrt{d/n}$, and a scheme estimating the same levels would pay $\sqrt{M/n}$ exactly as
stratifying does. The apparent saving is the assumption, not the method.

\paragraph{Claim that stratifying removes the need for the correction.} In the two-state model
the independent-data margin survives a pooled variance inflation of $223$ once the block is
stratified, which suggested the correction was redundant. That is an artefact of the model. A
two-state chain has one non-unit mode and it \emph{is} the state, so stratifying removes the
dependence by construction. On a continuous state the stratification must bin, and binning is
approximate: on an AR(1) log-volatility state at $\phi=0.98$ with two bins the independent margin
fails on $24\%$ of calibration draws against a nominal $5\%$, and the correction repairs it. The
failure decays as the binning refines and is gone by ten bins. The honest statement is conditional:
stratifying substitutes for the correction only when it is fine enough to resolve the state.

\paragraph{Explain the residual by clustering of stratum visits.} Members of a lagged-state
stratum arrive at return times whose mean is $M$ but whose median is $1$ at high persistence, with
over half of consecutive members adjacent in time. Since the autocorrelation across a gap is convex,
Jensen gives $\E[\phi^\tau] > \phi^{\E[\tau]}$, which looked like a clean account of why stratifying
buys less decorrelation than its sample-size reduction suggests. It is wrong. Kac's lemma forces
every equal-probability stratum to the same mean return time, so if visit clustering drove the
within-stratum variance then strata with matching gap statistics would have matching variance.
At $\phi=0.98$ with eight bins the mean gap varies by a factor $1.03$ across strata and
$\E[\phi^\tau]$ by $1.05$, while the within-stratum long-run variance of the coverage indicator
varies by orders of magnitude. What that variance tracks is $p_b(1-p_b)$ at each stratum's own
conditional coverage, which runs from $1.0000$ in the calmest to $0.4906$ in the most volatile, and
the ratio of long-run to Bernoulli variance sits near one in six of the eight strata. Almost no
serial dependence survives within a stratum. The correct object is the long-run covariance of the
coverage indicator, which is classical, and the extremal index is also wrong, being defined only for
sets of vanishing measure.

Each of the four tried to convert an exact description of the
corrector into a procedure that beats an existing one, and each was beaten by something simpler or
was measuring an artefact of the two-state cartoon. The description survives. The procedures do not.

\section{Interpretation}\label{sec:interp}

The corrected empirical quantile level is
\begin{equation}\label{eq:corrected-level}
  p + \underbrace{a_\ell}_{\text{fast-state protection}}
    + \underbrace{\delta_{n,\eta}}_{\text{dependent-calibration protection}} ,
\end{equation}
and in threshold units the first-order inflation is $(a_\ell+\delta_{n,\eta})/f(q_p)$. The two
costs are qualitatively different. The transient term $a_\ell$ is not an estimation error. It
is the real difference between the current conditional law and the invariant law. The term
$\delta_{n,\eta}$ is the price of demanding that the realized calibration sample be
trustworthy with probability $1-\eta$. A state-adaptive implementation need not pay the
worst-case $a_\ell$. It can estimate $b_{\ell,W_n}(q)=(P^\ell-P^{\ell+1})\chi_q(W_n)$ and build
a lower confidence envelope for the current conditional CDF. Theorem~\ref{thm:honest} is the
simplest version with a transparent proof.

\section{Combining with the local sharpness modulus}\label{sec:modulus}

The companion note \citep{cotton2026honest} shows that honest coverage at a fixed covariate
value $x_0$ forces inflation at the local nonparametric modulus $(n g(x_0))^{-s/(2s+d)}$, from
locally indistinguishable alternatives rather than any exchangeability phenomenon. The
dependent extension is schematic, and we label it as such. For a kernel $K_h$ the relevant
local object is the weighted long-run variance
\begin{equation}\label{eq:omegah}
  \Omega_h(x_0) = \sum_{j\in\mathbb{Z}} \Cov\Big(K_h(X_0-x_0)[\one\{R_0\le q_p(x_0)\}-p],\,
    K_h(X_j-x_0)[\one\{R_j\le q_p(x_0)\}-p]\Big),
\end{equation}
and the argument needs a local Markov mixing CLT of the form $h^d\Omega_h(x_0)\to
\omega_p(x_0)$.

\begin{remark}[Schematic bias--margin decomposition]\label{rem:schematic}
Under a fixed uniform spectral gap the effective local sample stays of order $n g(x_0)h^d$,
and the natural upper bound is, in threshold units,
\begin{equation}\label{eq:hierarchy-bound}
  T_{n,h}(x_0) - q_p(x_0) \lesssim
    C_q L h^{s}
    + \frac{\sqrt{\omega_p(x_0)\log(1/\eta)}}{f_{p,x_0}\sqrt{n g(x_0) h^{d}}}
    + \frac{a_{\ell,x_0}}{f_{p,x_0}}
    + \frac{R_{\mathrm{hom}}(\varepsilon) + \Delta_{\mathrm{dyn}}}{f_{p,x_0}} ,
\end{equation}
all terms on the threshold scale. Minimizing the first two over $h$ gives $h_\star\asymp
[\omega_p(x_0)/(n g(x_0)L^2)]^{1/(2s+d)}$ and local inflation $\asymp L^{d/(2s+d)}
[\omega_p(x_0)/(n g(x_0))]^{s/(2s+d)}$.
\end{remark}

So under a fixed spectral gap, dependence changes the information constant through the local
long-run variance but leaves the H\"older exponent unchanged. If the mixing time itself
diverges with $n$, $h$, or $\varepsilon$, then $\omega_p(x_0)$ may diverge and the rate can
change. Homogenization adds the transient and dependence terms and does not erase the local
modulus. The hierarchy is
\begin{equation}\label{eq:hierarchy}
  \begin{aligned}
    \text{local heterogeneity} &\;\to\; \text{nonparametric modulus},\\
    \text{calibration persistence} &\;\to\; \text{Poisson long-run variance},\\
    \text{current fast state} &\;\to\; \text{Poisson transient corrector},\\
    \text{approximating the fast--slow model} &\;\to\; \text{homogenization remainder}.
  \end{aligned}
\end{equation}
The first term is unavoidable over a broad H\"older class. The second and third can be
quantified and corrected under dynamical structure. The fourth must be included for the
guarantee to be honest.

\section{The limits of the parametric gain}\label{sec:whatbuys}

Homogenization does not defeat the sharpness lower bound. It partitions the conformal error
into unavoidable local ambiguity, correctable dynamical bias, and estimable dependence noise.
If the fast--slow dynamics live in a finite-dimensional parametric family, honest uncertainty
about the dynamical correction shrinks at the parametric rate $n^{-1/2}$ rather than the
nonparametric local rate $n^{-s/(2s+d)}$. That gain does not contradict the companion note. It
comes from replacing a broad H\"older class by a much smaller structured experiment. If the
transition law, the state-specific score distributions, or the scale field are again allowed
to vary nonparametrically with $x$, the local modulus returns, and if their smoothness is
unknown the no-adaptation obstruction returns with it.

\section{The corresponding skaters transform}\label{sec:skaters}

The correction is directly usable as a recalibration layer. For an existing distributional
forecast $F^0_{t+1}$, compute the Gaussianized probability integral transform $Z_t=
\Phi^{-1}(F^0_t(Y_t))$, and fit a predictable two-state model to the past $Z_t$ with
state-specific CDFs $K_L,K_H$ and filtered high-state probability $\omega_t$. With
$\omega_{t+1\mid t}=\pi_H+\lambda(\omega_t-\pi_H)$ the conditional CDF of the next Gaussianized
PIT is $K_{t+1\mid t}(z)=(1-\omega_{t+1\mid t})K_L(z)+\omega_{t+1\mid t}K_H(z)$, and the
dynamically corrected predictive CDF and $u$-quantile are
\begin{equation}\label{eq:corrected-cdf}
  F_{t+1}(y) = K_{t+1\mid t}\big[\Phi^{-1}(F^0_{t+1}(y))\big],
  \qquad
  \widehat Q_{t+1}(u) = (F^0_{t+1})^{-1}\big[\Phi(K^{-1}_{t+1\mid t}(u))\big] .
\end{equation}

This is a predictable dynamic PIT recalibration, not by itself a finite-sample conformal
procedure. The fence-post map estimates a pooled PIT law, and \eqref{eq:corrected-cdf}
estimates the conditional PIT law after accounting for the current residual regime. To connect
it to Theorem~\ref{thm:honest} the implementation needs a predictable or separately estimated
regime model, an uncertainty envelope covering parameter and filtering error, a
probability-level margin, and a treatment of ties or discrete forecast distributions. For an
honest upper quantile, let $\mathcal{C}_t$ be a confidence set for the dynamic parameters and
add a homogenization or misspecification radius $r_t$,
\begin{equation}\label{eq:honest-skater}
  z^{\mathrm{hon}}_{t+1,p} = \sup_{\theta\in\mathcal{C}_t} K^{-1}_{\theta,t+1\mid t}(p+r_t),
  \qquad
  y^{\mathrm{hon}}_{t+1,p} = (F^0_{t+1})^{-1}\big[\Phi(z^{\mathrm{hon}}_{t+1,p})\big] .
\end{equation}
Because distributional conformal prediction for Markov processes already composes a
transition-CDF PIT \citep{dai2024markov}, the novelty here is the Poisson state correction and
its honesty envelope, not the conditional PIT composition itself.

\paragraph{Empirical protocol.} A comparison should include the base forecast, the marginal
fence-post correction, an AR correction, the two-state homogenized correction, and the honest
two-state correction. The decisive diagnostics are PIT calibration conditional on the filtered
state, coverage conditional on recent $Z_t$ and volatility, the lower-tail probability of
realized coverage, proper scores such as log score and CRPS, and the extra interval width the
honest envelope demands. Aggregate PIT uniformity alone is insufficient, since the two-state
example shows pooled $90\%$ calibration concealing conditional coverage from $74\%$ to $94\%$.

\section{Reproducibility}\label{sec:repro}

Two short scripts accompany the note. \texttt{verify\_decomposition.py} checks, on random
finite chains, the corrector identity \eqref{eq:corrector-obs} (maximum error $5.6\times
10^{-16}$), the identity $\psi_p=\chi_{q_p}$ of \eqref{eq:psi-is-chi} (exactly zero), the
reversible spectral identities of Proposition~\ref{prop:spectral} (error $\sim10^{-16}$), and
the generator limit of Proposition~\ref{prop:homog} (error decaying linearly in $\delta$).
\texttt{two\_state\_illustration.py} reproduces the $74\%$/$94\%$ conditional coverages, the
long-run variance $0.17$, and the effective sample fraction $0.53$, and confirms $\sigma_p^2$
by direct simulation of the sampled chain. \texttt{garch.py} produces the table of Section~\ref{sec:garch}.
\texttt{mechanism.py} runs both margins in both schemes across the persistence range, and
\texttt{jensen\_audit.py} is the test that retracted the clustering explanation.
\texttt{compare\_arms.py} produces the threshold comparison of Section~\ref{sec:failed}, and
\texttt{stress.py} is the attempt to break the independent margin in the stratified scheme, which
failed at a variance inflation of $223$. \texttt{EXPERIMENTS.md} records the full sequence,
including a variance bug in an earlier draft: a potential computed on the regime alone drops the
emission term and understates $V$ by an order of magnitude at low persistence.

\section{What would make the reading useful}

The averaged limit is valid when the fast variable is fast, so the quantity a practitioner needs is
the volatility relaxation time against the forecast horizon. That ratio is estimable from the same
data used for calibration, and the GARCH table says roughly what to expect from it: a half-life of
three days costs six points of spread in conditional coverage, seventeen days costs sixteen, and
seventy days costs twenty-four. A diagnostic reporting the ratio alongside the coverage rate would
tell a user whether the pooled interval can be read conditionally, which is the reading most users
take from it anyway.

One question is open and matters for how much of this is new. Theorem~\ref{thm:corrector} is an
identity about occupation functionals of a Markov chain, and identities of that shape are old. We
have not searched that literature.

\bibliographystyle{plainnat}
\begin{thebibliography}{9}
\bibitem[Barber et al.(2023)]{barber2023beyond} Barber, R.~F., Cand\`es, E.~J., Ramdas, A.,
  and Tibshirani, R.~J. (2023). Conformal prediction beyond exchangeability. \emph{Annals of
  Statistics} 51(2), 816--845.
\bibitem[Chernozhukov et al.(2021)]{chernozhukov2021distributional} Chernozhukov, V.,
  W\"uthrich, K., and Zhu, Y. (2021). Distributional conformal prediction. \emph{Proceedings
  of the National Academy of Sciences} 118(48).
\bibitem[Cotton(2026)]{cotton2026honest} Cotton, P. (2026). Locally honest coverage forces
  nonparametric inflation: a sharpness lower bound for prediction sets. Working paper.
\bibitem[Bayley and Hammersley(1946)]{bayley1946effective} Bayley, G.~V. and Hammersley, J.~M.
  (1946). The ``effective'' number of independent observations in an autocorrelated time series.
  \emph{Journal of the Royal Statistical Society, Supplement} 8(2), 184--197.
\bibitem[Dai et al.(2026)]{dai2024markov} Dai, Wu, and Politis (2026). Distributional
  conformal prediction for Markov processes. Preprint.
\bibitem[Freedman(1975)]{freedman1975} Freedman, D.~A. (1975). On tail probabilities for
  martingales. \emph{Annals of Probability} 3(1), 100--118.
\bibitem[Glynn and Henderson(2002)]{glynn2002quantile} Glynn, P.~W. and Henderson, S.~G.
  (2002). Quantile estimation for dependent data. In \emph{Analysis of Simulation Output}.
\bibitem[Meyn and Tweedie(2009)]{meyntweedie2009} Meyn, S. and Tweedie, R.~L. (2009).
  \emph{Markov Chains and Stochastic Stability}, 2nd ed. Cambridge University Press.
\bibitem[Oliveira et al.(2024)]{oliveira2024split} Oliveira, R.~I., Orenstein, P., Ramos, T.,
  and Romano, J.~V. (2024). Split conformal prediction and non-exchangeable data. \emph{Journal
  of Machine Learning Research} 25.
\bibitem[Ramos et al.(2026)]{ramos2024transported} Ramos, T., Graziadei, H., and Cabezas, L.
  (2026). Conformal prediction via transported beta laws. Preprint.
\bibitem[Zheng and Proutiere(2024)]{zheng2024markovian} Zheng, X. and Proutiere, A. (2024).
  Conformal prediction under Markovian data. Preprint.
\end{thebibliography}

\end{document}
